% sidenote_type=symbol（傍注，4.2.7）: pdfTeX検証
% 過去のpdfTeX不具合(b)の回帰フィクスチャ：symbol型\sidenoteが
% 「Undefined control sequence」（\CJK@punctchar内部エラー）．
% 既定マーク ＊（U+FF0A，CJK utf8 のアクティブ文字）を素の \edef で
% \@thefnmark に入れており，アクティブ文字の展開連鎖が \edef 下で
% \CJK@punctchar の内部定義（\def\CJK@tempa##1...）に踏み込んで壊れた．
% \protected@edef 化で文字トークンを保持するよう修正済み（pTeX/upTeX/
% LuaTeXでは和文文字は展開不能なので従来と同値）．
\documentclass[pdflatex,sidenote_length=20zw]{jlreq}
\jlreqsetup{sidenote_type=symbol}
\begin{document}
% --- sidenote_type=symbol ---
記号型の傍注では\sidenote{強調されている語句}{その語句についての説明注。
見出し付きで欄外に出る。}という形になる。

分離型も：\sidenotemark{別の語句}文中に置き，\sidenotetext{別の語句}{分離した記号型の傍注本文。}
注はあとで与える。

% --- 共存する脚注・後注・単体 marginpar ---
脚注\footnote{脚注本文。}と後注\endnote{後注本文。}，
単体の\marginpar{欄外テスト}も通す。

\theendnotes
\end{document}
